% SPDX-License-Identifier: Apache-2.0
%
% The TxHDL cheat sheet: two pages, landscape, meant to be printed on
% the two sides of one sheet. The front walks one small unit from its
% source to its netlist and its waveform; the back is the vocabulary.
% Each side stands on its own. Built by //docs:cheatsheet, and as one
% PNG per side by //docs:cheatsheet_png.
\documentclass[10pt]{article}
\usepackage[paperwidth=11in,paperheight=8.5in,margin=0.4in,top=0.35in,bottom=0.35in]{geometry}
\usepackage{multicol}
\usepackage{array}
\newcommand{\listingsize}{\fontsize{6.2}{7.2}\selectfont}
% SPDX-License-Identifier: Apache-2.0
%
% The house preamble, shared by every document in this directory. Loaded
% after \documentclass, before \title. Keeping it in one file is what
% stops two articles drifting apart in font, listing style or figure
% style.
% Computer Modern. IEEEtran selects Times (ptm) by default, so the face has
% to be chosen explicitly.
%
% Latin Modern is the Computer Modern variety used here, and the choice is
% forced rather than aesthetic. Under T1 encoding, plain `cmr` has no Type 1
% outlines unless cm-super is installed, and it is not in the pinned
% distribution, so pdflatex falls back to EC bitmaps and every glyph in the
% PDF becomes a Type 3 font. That was measured: the first build produced a
% document with 10 Type 3 fonts and no Type 1. Latin Modern is Computer
% Modern redrawn with a full T1 repertoire and real .pfb outlines, so the
% same design comes out scalable.
\usepackage[T1]{fontenc}
\usepackage[utf8]{inputenc}
\usepackage{lmodern}
% microtype is not in the pinned TeX distribution. emergencystretch alone
% keeps the two-column measure from producing overfull lines.
\setlength{\emergencystretch}{3em}

\usepackage{amsmath}
\usepackage{amssymb}
\usepackage{amsthm}
\usepackage{booktabs}
\usepackage{array}
% enumitem is not in the pinned TeX distribution; the list spacing below
% does what \setlist would have done.
\usepackage{float}
\usepackage{xcolor}
\usepackage{listings}
\usepackage{tikz}
\usetikzlibrary{arrows.meta,positioning,fit,backgrounds,calc}
\usepackage[hidelinks,bookmarks=true]{hyperref}

% Numbered, definition-style box for a rule the reader is meant to apply.
\theoremstyle{definition}
\newtheorem{rulebox}{Rule}
\newtheorem{example}{Example}

% Unnumbered asides.
\theoremstyle{remark}
\newtheorem*{tip}{Tip}
\newtheorem*{pitfall}{Pitfall}

% Tighter lists, in place of enumitem's \setlist.
\setlength{\topsep}{2pt}
\setlength{\itemsep}{1pt}
\setlength{\parsep}{0pt}

% Code in the running text. The typewriter face under T1 ligates `<<`
% and `>>` into guillemets, so a shift printed as \code{<<} came out as
% a French quotation mark (issue 571). microtype's \DisableLigatures is
% not in the pinned distribution, but the primitive it rests on is
% pdfTeX's own: \pdfnoligatures strips every ligature from the font it
% is given, and the current font inside \texttt is the typewriter face
% at the size in use, so each size is stripped the first time \code
% sets it. Code wants no ligature at all: two quotes are two quotes.
\newcommand{\code}[1]{\texttt{\ifdefined\pdfnoligatures\pdfnoligatures\font\fi #1}}
\newcommand{\term}[1]{\emph{#1}}

% Syntax highlighting. Four roles, distinguishable in colour and still
% distinguishable in greyscale, because an IEEE article gets printed: the
% keyword is the only bold role, the comment is the only italic one, and the
% two remaining roles differ in lightness as well as in hue.
\definecolor{kwblue}{RGB}{20,60,150}
\definecolor{cmtgreen}{RGB}{90,120,90}
\definecolor{strred}{RGB}{150,50,40}
\definecolor{numgray}{gray}{0.45}
\definecolor{framegray}{gray}{0.70}
\definecolor{bgshade}{gray}{0.97}
% A darker fill for figure cells. The listing background has to stay very
% light to keep code readable; a figure cell has to be clearly shaded or the
% distinction it draws is invisible in print.
\definecolor{cellshade}{gray}{0.90}

% The listing size is the document's choice, because it depends on the
% column: a two-column article sets `\listingsize` to 6pt and keeps its
% listings within 70 columns, or to 5.6pt when it includes source that
% rustfmt sets at 80, which is what fits a column's frame (issue 1061);
% a one-column document keeps the body size and includes source files
% at 80. `//docs:frame_test` checks every listing line against its frame. Lines are never broken; a line that does not fit
% overflows the frame, which is visible, rather than wrapping, which is
% not.
\providecommand{\listingsize}{\scriptsize}

% `'` is deliberately not a string delimiter: the language uses it for loop
% labels, as Rust does, so treating it as a quote would swallow the rest of
% the line.
\lstdefinestyle{txhdl}{
  basicstyle=\ttfamily\listingsize,
  keywordstyle=\bfseries\color{kwblue},
  commentstyle=\itshape\color{cmtgreen},
  stringstyle=\color{strred},
  numberstyle=\tiny\color{numgray},
  morekeywords={mod,use,pub,const,type,struct,enum,trait,impl,fn,async,let,mut,
    match,if,else,loop,while,for,in,break,continue,return,as,await,
    module,bus,channel,transaction,protocol,pipeline,flow,seq,config,
    port,stage,pipe,instance,connect,bind,tag,domain,blackbox,foreign,
    property,role,signal,handler,true,false,
    sequence,interface,modport,implements,package,import,var,wire,func,proc,
    case,default,next,fork,branch,join,state,fsm},
  morecomment=[l]{//},
  morestring=[b]",
  showstringspaces=false,
  columns=fullflexible,
  keepspaces=true,
  breaklines=false,
  % Framed, so a listing is bounded rather than floating in the column.
  frame=single,
  rulecolor=\color{framegray},
  framesep=3pt,
  backgroundcolor=\color{bgshade},
  xleftmargin=3pt,
  xrightmargin=3pt,
  aboveskip=6pt,
  belowskip=6pt,
  % Every listing is numbered and captioned, so it can be referred to by
  % number. `caption` without `float` numbers the listing and leaves it
  % where it was written, which is what a listing in running prose needs.
  captionpos=b,
  abovecaptionskip=2pt,
  belowcaptionskip=6pt,
}

% Figures drawn as text. Framed the same way, and with the highlighting
% switched off, because a box-and-arrow diagram is not code and half its
% words turning blue reads as an accident. Setting a style adds keys rather
% than clearing the ones already set, so the three roles are neutralised by
% hand rather than by leaving them out. Program output is printed in this
% style too, and a netlist's assignment can run past the frame, so a line
% that does not fit is broken and the rest indented; source never is.
\lstdefinestyle{ascii}{
  basicstyle=\ttfamily\listingsize,
  keywordstyle=\ttfamily,
  commentstyle=\ttfamily,
  stringstyle=\ttfamily,
  columns=fullflexible,
  keepspaces=true,
  breaklines=true,
  breakindent=2em,
  frame=single,
  rulecolor=\color{framegray},
  framesep=3pt,
  backgroundcolor=\color{bgshade},
  xleftmargin=3pt,
  xrightmargin=3pt,
  aboveskip=6pt,
  belowskip=6pt,
}
\lstset{style=txhdl}

% House TikZ styles. `shorten` lives on the arrow styles rather than on each
% arrow, so every arrow in the document gets the gap, including ones added
% later. TikZ clips a path at a node's border, so without it an arrowhead
% butts against the box with no gap at all. Keep `node distance` well above
% twice the shorten, or the arrow shrinks to nothing.
\tikzset{
  box/.style={draw,rounded corners=2pt,align=center,inner sep=4pt,
              font=\footnotesize},
  boxf/.style={box,fill=bgshade},
  lbl/.style={font=\scriptsize\itshape,align=center},
  fl/.style={-{Stealth[length=4pt]},shorten >=2.5pt,shorten <=2.5pt},
  fld/.style={fl,dashed},
}


% The contents list: the class gives a section's number 2.75em, which
% a Roman numeral past thirty overruns onto the title, so the box is
% wider here, and a subsection's entry starts where a title does.
\makeatletter
\def\l@section#1#2{\addpenalty{\@secpenalty}\addvspace{1.0em plus 1pt}%
    \@tempdima 5.75em \begingroup \parindent \z@ \rightskip \@pnumwidth%
    \parfillskip-\@pnumwidth {\bfseries\leavevmode #1}\hfil\hbox to\@pnumwidth{\hss #2}\par%
    \endgroup}
\def\l@subsection{\@dottedtocline{2}{5.75em}{5em}}
\def\l@subsubsection{\@dottedtocline{3}{10.75em}{5.5em}}
\makeatother

\pagestyle{empty}
\setlength{\columnsep}{16pt}
\setlength{\parindent}{0pt}
\setlength{\parskip}{2pt}
\definecolor{bar}{RGB}{20,60,150}
\definecolor{tint}{RGB}{232,238,250}

% A section head: a coloured bar.
% A heading never sits alone at the foot of a column: when less than
% five lines are left, the column breaks before it.
\newcommand{\keep}{\par\vskip 0pt plus 5\baselineskip\penalty-200\vskip 0pt plus -5\baselineskip}
\newcommand{\head}[1]{\keep\vspace{4pt}\noindent\colorbox{bar}{\parbox{\dimexpr\linewidth-2\fboxsep}{\color{white}\bfseries\small #1}}\vspace{2pt}\par}

% The banner across the top of a side.
\newcommand{\banner}[2]{\noindent\colorbox{bar}{\parbox{\dimexpr\textwidth-2\fboxsep}{\color{white}
{\LARGE\bfseries #1}\hfill{\small #2}}}\vspace{4pt}\par}

% A two-column table of code and what it means, a column wide.
\newenvironment{vocab}{\par\footnotesize\setlength{\tabcolsep}{3pt}\begin{tabular}{@{}>{\raggedright\arraybackslash}p{0.53\linewidth}>{\raggedright\arraybackslash}p{0.44\linewidth}@{}}}{\end{tabular}\par}

\begin{document}
\small

% ---------------------------------------------------------------- front
\banner{TxHDL Cheat Sheet}{Hardware in Rust: verified module walkthrough on front; syntax and vocabulary on back}

\begin{multicols}{3}

\head{The idea}
TxHDL is hardware in Rust. Where Rust has a construct, TxHDL uses it
directly: a transaction is a \code{struct}, a unit is a \code{struct}
implementing \code{Unit}, a pipeline is an \code{async fn}, a clock is an
affine type. The compiler enforces design rules: two drivers on a wire is a
move error; crossing clock domains without synchronizers is a type mismatch;
cycles in combinational logic cannot compile without \code{.await}.

\head{The one rule about time}
A process \textbf{waits for an edge, then reads, then drives}. Every
\code{.await} marks a cycle boundary. After the edge, registers reflect
sampled state, driven registers update at cycle conclusion, and driven
outputs update combinationally. This ordering preserves synchronous
RTL semantics while allowing reads and drives in one step.

Waits: \code{C::rising()}, \code{C::falling()}, and
\code{until(C::rising, || cond)} for conditional edges. Concurrent edge
evaluation uses \code{parallel!(a, b).await} (simulation only). Time
advances in half-cycle ticks (two ticks per clock cycle).

\head{Build it, run it, look at it}
Prerequisites are \code{bazelisk} and Git. Toolchains, simulators, and
\LaTeX{} are hermetic.
\begin{lstlisting}[style=ascii,frame=none,backgroundcolor=\color{tint},xleftmargin=2pt]
bazel build //...                 # everything
bazel test  //...                 # and check it
bazel run //lib/examples:ex_cheat # the netlist
TXHDL_FST=$PWD/c.fst bazel run \
  //lib/examples:ex_cheat         # and a waveform
bazel build //docs/...            # every document
\end{lstlisting}

\head{Where to read more}
\code{//docs:tutorial} walks through in nine steps. \code{//docs:examples}
provides compiled reference designs. \code{//docs:zero} bootstraps from an
empty workspace. \code{//docs:cover} catalogues every document.

\head{Adding your own}
Copy \code{lib/examples/ex\_cheat.rs} to \code{ex\_\emph{name}.rs}, register
it in \code{lib/examples/BUILD.bazel}, add a \code{waveform()} target, and
document its section in \code{//docs:examples}. Under the triad rule: runtime,
example, and documentation advance together.

\columnbreak

\head{A whole unit: a counter}
Counts while \code{enable} is asserted and raises \code{tick} on terminal
seven. Top to bottom: the struct declares state, \code{impl} defines
synchronous behavior, and function signatures expose I/O ports.
\lstinputlisting[rangebeginprefix=//\ begin\{,rangebeginsuffix=\},rangeendprefix=//\ end\{,rangeendsuffix=\},includerangemarker=false,linerange=unit-unit,frame=none,backgroundcolor=\color{tint},xleftmargin=2pt]{lib/examples/ex_cheat.rs}
\code{\#[lower]} derives synthesizable netlists directly. Unannotated
modules simulate if ports are declared as \code{Unit<In, Out>}.

\head{What it turns into}
Synthesised netlist from \code{Counter::verilog("counter")}. Synchronous
reset \code{rst} is bound implicitly. Bazel derives equivalent VHDL entities
and co-simulates both against the Rust trace, cycle by cycle.
\lstinputlisting[style=ascii,frame=none,backgroundcolor=\color{tint},xleftmargin=2pt]{out_cheat.txt}

\columnbreak

\head{And what it did}
Rendered directly from the simulation trace. Signal \code{enable} deasserts
at cycle 10, \code{n} counts while high, and \code{tick} pulses on the edge
following seven.
\begin{center}
\input{cheat_timing_color.tex}
\end{center}

\head{The shape every unit has}
\begin{itemize}\setlength{\itemsep}{0pt}\setlength{\leftmargin}{8pt}
\item One \code{loop}, opening with an edge wait.
\item Sample inputs via \code{get}, read registers by name.
\item Combinational \code{let}s, each a named wire in the netlist.
\item Single-assignment drives: \code{with!} for registers, \code{set}
  for outputs, \code{send} for channels under \code{ready}.
\end{itemize}
Single assignment and branch-free drives ensure deterministic lowering.

\head{How the netlist is checked}
For every lowered design, Bazel captures the golden Rust trace, derives VHDL
and Verilog netlists, and replays input vectors, checking registers and pins
at each tick: VHDL under \code{nvc}, Verilog under Verilator. Adding
\code{lowered = (entity, unit)} to \code{waveform()} enables both tests.

\head{When something goes wrong}
Procedural macros reject non-lowerable syntax at compile time with targeted
diagnostics. In co-simulation, tests fail at the first diverging tick. The
generated waveform isolates root causes immediately.

\vfill
\InputIfFileExists{buildstamp.tex}{}{}
\providecommand{\buildstamp}{Draft build (outside the Bazel flow).}
{\footnotesize\buildstamp\ Written with a large language model,
Claude, as an assistant, as every commit says. Turn over for the
vocabulary.}

\end{multicols}

% ---------------------------------------------------------------- back
\newpage
\banner{TxHDL Vocabulary}{Complete language syntax, hardware primitives, and lowering semantics}
\footnotesize
% The back holds the whole vocabulary, so its tables are a size smaller.
\renewenvironment{vocab}{\par\scriptsize\setlength{\tabcolsep}{3pt}\begin{tabular}{@{}>{\raggedright\arraybackslash}p{0.55\linewidth}>{\raggedright\arraybackslash}p{0.42\linewidth}@{}}}{\end{tabular}\par}

\begin{multicols}{3}

\head{Values and state}
\begin{vocab}
\code{Bit}, \code{U<N>}, \code{I<N>} & scalar bits and words; operator overloads \\
\code{\#[derive(Value)]} & custom enum/struct value, named in traces \\
\code{\#[derive(Transaction)]} & structured type for channel transport \\
\code{Reg<T, C>} & register on clock \code{C} (default omitted); direct read \\
\code{Wire<T>} & combinational wire as struct field for tracing \\
\code{Mem<T, N>} & synchronous RAM; \code{read(a)}, \code{with!(.. at(a) <= d)} \\
\code{slice}, \code{concat}, \code{sext}, \code{zext} & bit extraction, concat, sign/zero extension \\
\code{mul::<M>()} & integer multiplication yielding $M$-bit product \\
\code{U<256, 2>} & wide words ($N$-bit vector across $L$ 128-bit limbs) \\
\end{vocab}

\head{Driving}
\begin{vocab}
\code{with!(self <= \{ c ? n: v \})} & conditional register update \\
\code{when!(c => self \{ .. \})} & condition-first synchronous update \\
\code{case!(v => \{ p => \{..\} \})} & priority pattern register update \\
\code{if c \{ .. \}} & conditional block around drives \\
\code{out.set(v)} & drive combinational pin in current cycle \\
\end{vocab}

\head{Choosing a value}
\begin{vocab}
\code{mux(c, a, b)} & 2-to-1 combinational multiplexer \\
\code{select!(v => \{ p => e, .. \})} & pattern-matching multiplexer \\
\code{if c \{ a \} else \{ b \}} & conditional expression lowered as \code{mux} \\
\code{match v \{ .. \}} & pattern match lowered as \code{select!} \\
\code{0x30..=0x39}, \code{(lo..=hi).contains(\&x)} & bit-range patterns or guards \\
\end{vocab}
Hardware multiplexes rather than branches: both expressions evaluate concurrently and condition selects. Syntax mirrors circuit topology.

\head{Wires and channels}
\begin{vocab}
\code{signal()} $\to$ \code{(Out, In)} & combinational wire; affine \code{Out} prevents multi-driver conflicts \\
\code{chan()} $\to$ \code{(Tx, Rx)} & channel with 2-element elastic FIFO \\
\code{tx.send(v)}, \code{tx.ready()} & non-blocking enqueue; test buffer room \\
\code{tx.put(|| v).await} & blocking send; suspends until space opens \\
\code{rx.peek()}, \code{rx.recv()} & inspect head; consume and dequeue \\
\code{rx.recv\_if(c)}, \code{rx.take()} & conditional dequeue; unconditional dequeue \\
\code{pad()} & bidirectional tri-state I/O pad \\
\end{vocab}
In netlists, channels map to data, \code{valid}, and \code{ready}. An \code{until} wait guards subsequent logic.

\head{Ports with names}
Scalar ports suit simple counters. Complex interfaces bundle signals into structs to eliminate pin-swapping bugs:
\begin{vocab}
\code{run(\&mut self, i: InP, o: OutP)} & group ports across structured types \\
\code{DiffIn \{ minuend, subtrahend \}} & destructure port fields in signature \\
\code{\#[derive(Ports)]} & derive macro for external port structs \\
a side \code{links: Links} & flattened hierarchy (\code{links\_left\_aw\_valid}) \\
\code{side: [Rx<T>; N]} & indexed port array (\code{side\_0} onward) \\
\code{interface! \{ Ops \{..\} role S \{..\} \}} & declare protocol once, derive role structs \\
\end{vocab}
Lowering flattens port structs into discrete pins; destructuring retains clean net names. AXI-Lite modules consume \code{LitePort}.

\head{Units of units}
Parent modules hold children as fields, wire ports in \code{run}, and evaluate via \code{join2} (nestable). Lowering derives module hierarchies.
\begin{vocab}
\code{let (h, p) = link::<LitePort<..>>()} & complete bidirectional link in one line \\
\code{child.run((x, tie(v)), y)} & tie input pin permanently to constant \\
\code{child.run((bus, ..), ..)} & forward structured port bundle intact \\
\end{vocab}

\head{Clocks}
\begin{vocab}
\code{struct Slow; impl Clock for Slow} & define clock: name, period, phase in ticks \\
\code{Reg<T, Slow>}, \code{In<T, Slow>} & register state and ports on clock \\
\code{Slow::rising().await} & process suspended awaiting active edge \\
\code{ChanCdc} & asynchronous dual-clock CDC channel \\
\end{vocab}
Multi-clock designs compose children across domains under an asynchronous top without edge waits. \code{DefaultClock} is omitted unless explicitly required.

\head{Reset}
Reset is implicit. Every synthesised module binds an active-high \code{rst} net that clears registers and empties channels. Custom reset logic declares \code{rst: In<Bit>}, tapping the global net directly.

\head{Other people's hardware}
External Verilog/VHDL modules integrate as child units: provide a simulation model in Rust and implement \code{Lower} via \code{netlist::foreign(..)}. Conversely, \code{verilog\_unit()} and \code{vhdl\_unit()} wrap foreign HDL modules into callable TxHDL units.

\head{What the lowering reads}
A single loop of an edge wait, reads, combinational bindings, and drives lowers to a synchronous sequential block. Multiple wait points synthesize as finite state machines. Loops of \code{for i in 0..N} across port arrays unroll statically. Supported operators: \code{+ - * \% \& | \^{} ! <{}< >{}>}, unary \code{-}, comparisons, and \code{as} casts. Functions marked \code{\#[lower]} inline transparently. Loops containing waits synthesize as multi-cycle counters bounded by registers or constants.

\head{Names}
HDL keyword collisions (e.g., \code{next}, \code{begin}, \code{signal}) resolve by appending \code{\_rw} in netlists and testbenches. Custom signal names are assigned via \code{\#[rename("...")]}.

\head{Register maps and checks}
\begin{vocab}
\code{regmap! \{ m (m\_read, m\_we[, m\_re]), 2: [..] \}} & single-source map: decoder, C header, datasheet \\
\code{(1, ctrl, rw, "..", [(run, 0, 1, ..)])} & register bit fields, access mode, description \\
\code{check!(c, "why")} & safety assertion; synthesizes as \code{assert} \\
\code{assume!(c, "why")}, \code{cover!(c, "why")} & formal constraint; formal reachability target \\
\end{vocab}

\head{Pipelines and builds}
\begin{vocab}
\code{\#[pipeline(..)] async fn} & statically scheduled pipeline; \code{.await} marks stages \\
\code{pipeline::drive(f, rx, tx)} & elastic runner; stalls when \code{tx} fills \\
\code{config! \{ Fpga: C for Top<Fpga> \{..\} \}} & declarative build configuration bundle \\
\end{vocab}

\head{Watching and testing}
\begin{vocab}
\code{\#[derive(Trace)]} & expose struct fields in simulation waveforms \\
\code{Wave::from\_env()} & waveform logger controlled by env vars \\
\code{w.add("name", \&x)}, \code{start()}, \code{stop()} & register signals; control capture \\
\code{Running::new(u.run(..))} & initialize unit execution state machine \\
\code{sim.cycle()} & advance one cycle; inspect registered state \\
\end{vocab}
Units execute in standard \code{\#[test]} harnesses, each maintaining isolated simulation time. Assertions verify registered states after cycle advancement.

\end{multicols}
\end{document}
